\section{Introduction}
\label{sec:intro}

Large language models (LLMs) have made remarkable progress under the autoregressive (AR) paradigm~\citep{brown2020language,touvron2023llama}, but the fixed left-to-right generation order imposes an inherent serial dependency on inference. Diffusion language models (DLMs) \citep{austin2021structured,sahoo2024simple,nie2026large,arriola2025block} offer a new way to break this constraint. Cola DLM~\citep{guo2026continuous} further decomposes text generation into semantic organization in a continuous latent space and text realization in a discrete space; through latent prior modeling and conditional decoding, it provides a viable path toward hierarchical language generation beyond direct token modeling.

However, a structural division between continuous and discrete spaces does not automatically yield an effective division of labor. The key questions are what information the latent space should compress and how the discrete generator should use it. Shortening the latent sequence is not equivalent to obtaining a compact semantic representation, and a reconstruction-only objective does not guarantee that the latents can predict missing or future text~\citep{meshchaninov2026cosmos,jiang2026textldm,liang2026aurora}. Meanwhile, the decoder must not only faithfully realize the latent semantics, but also generate the remaining content from the text context when the latent condition is incomplete or imprecise~\citep{uziel2026crocodil}. Semantic compression and discrete generation therefore cannot be designed independently: the information the latents should retain must be determined jointly with the information the decoder can complete.

This interdependence further raises the question of how continuous and discrete representations should be learned jointly. Fixing the encoder in advance determines the information carried by the latents before the discrete generator is fully learned, limiting the continual adaptation of the representation to decoding needs. At the same time, letting the latent-prior fitting objective directly shape the encoder may push the latents toward directions that are easier to fit but not necessarily more beneficial for discrete generation. Hierarchical language modeling thus requires not only learning continuous representations and a discrete generator, but also specifying which objectives should jointly shape the representation and which learning processes should remain decoupled.

To address these issues, we propose \method, which jointly studies continuous semantic compression, the discrete reconstruction--generation balance, and their co-adaptation within a unified hierarchical language modeling framework. Its core idea is that, under limited latent capacity, the actual demands of text reconstruction and generation jointly determine what the continuous representation should retain, while the discrete generator handles the residual uncertainty that the latents do not explain. Specifically, \method dynamically segments text into semantic spans, controls the compression rate with an auxiliary-loss-free negative-feedback mechanism, and combines reconstruction and generation supervision to learn compact and predictive latent representations. On the discrete side, a block-causal masked diffusion model (MDM) decoder \citep{sahoo2024simple,arriola2025block} is jointly trained under clean, noisy, and dropped latent conditions, balancing semantic reconstruction with context-based autonomous generation, and classifier-free guidance (CFG) \citep{ho2022classifier} adjusts this balance at inference.

During training, the token-prediction loss is backpropagated through the latents to the encoder, so that the continuous representation and discrete generation co-adapt; the Diffusion Transformer (DiT) \citep{peebles2023scalable} learns the current latent distribution via Flow Matching \citep{lipman2023flow} without backpropagating gradients to the encoder, thereby decoupling prior fitting from representation learning. At inference, the model re-encodes the generated text block by block under the fixed span layout, keeping the latent history consistent with the actual text context and aligning the two spaces throughout generation.

Building on these ideas, we systematically study co-adaptive hierarchical language generation across continuous and discrete spaces from three perspectives: probabilistic modeling, representation learning, and generation performance. Our main contributions are as follows:
\begin{itemize}
    \item \textbf{We propose \method, a co-adaptive hierarchical language modeling framework.} From the perspective of conditional information decomposition, we analyze semantic compression, the reconstruction--generation balance, and continuous--discrete co-adaptation in a unified manner, clarifying the relation between the information retained by the latents and the residual left to discrete generation.
    \item \textbf{We develop a generation-oriented joint learning mechanism for continuous and discrete representations.} It combines dynamic semantic compression, generation-oriented supervision, and training under latent conditions of varying reliability, and coordinates representation learning with prior fitting through gradient decoupling, so that compact semantic representations and discrete generation capability co-evolve.
    \item \textbf{We validate the core designs and their scaling potential through systematic experiments.} Controlled five-task scaling to about 210B tokens and a separate 346.03B-source-token pretraining run show competitive, unsaturated scaling. The latter reaches a 60.10 average over eight tasks and surpasses Llama-3-8B on HellaSwag and StoryCloze despite over $43{\times}$ fewer reported tokens, demonstrating decoder capability transfer and compatibility with autoregressive and masked-diffusion token-space generators.
\end{itemize}
